
\section{Self-improvement}
\label{sec:self-improvement}

Throughout this section, $(\psi, A, B, L)$ will denote a fixed $(\eps, \delta, \gamma)$-good symmetric strategy for the $(m,q,d)$ low individual degree test.
The majority of this section will be devoted to proving \Cref{lem:self-improvement-helper} below,
which is a slightly weaker form of \Cref{thm:self-improvement-in-induction-section}.
The key difference is that the measurement~$H$ it outputs is allowed to be non-projective,
rather than the projective measurement given by \Cref{thm:self-improvement-in-induction-section}.
Having proven this, we can apply \Cref{thm:orthonormalization}
to produce a projective measurement;
this is done in \Cref{sec:self-improvement-projective} below,
completing the proof of \Cref{thm:self-improvement-in-induction-section}.

We now highlight other differences between \Cref{lem:self-improvement-helper}
and \Cref{thm:self-improvement-in-induction-section}.
Since~$H$ is non-projective, 
we have stated its strong self-consistency in \Cref{item:self-improvement-self}
in terms of \Cref{def:strong-self-consistency}; see \Cref{prop:two-notions-of-self-consistency}
for a proof that these conditions are equivalent for projective sub-measurements.
The other key difference is that the boundedness condition is modified slightly in \Cref{item:self-improvement-boundedness}.
Finally, the error $\zeta$ is substantially smaller.

\begin{lemma}[Self-improvement with non-projective output]\label{lem:self-improvement-helper}
Let $G \in \polymeas{m}{q}{d}$ be a measurement with the following property:
\begin{enumerate}

  \item (Consistency with~$A$): \label{item:self-improvement-G-consistency} On average over $\bu \sim \F_q^{m}$,
	  \begin{equation*}
	  A^{u}_a \otimes I \simeq_{\nu} I \otimes G_{[g(u)=a]}.
	  \end{equation*}
	  \end{enumerate}
Let
\begin{equation*}
\zeta = 100m\cdot \Big(\eps^{1/2} + \delta^{1/2} + (d/q)^{1/2}\Big).
\end{equation*}
Then there exists $H \in \polysub{m}{q}{d}$ with the following properties:
\begin{enumerate}
    \item(Completeness):  \label{item:self-improvement-completeness} If $H =  \sum_h H_h$, then
    \begin{equation*}
    \bra{\psi} H \otimes I \ket{\psi} \geq (1-\nu)-\zeta.
    \end{equation*}
    \item(Consistency with~$A$):\label{item:self-improvement-A-consistency} On average over $\bu \sim \F_q^m$,
        \begin{equation*}
            A^u_a \otimes I \simeq_{\zeta} I \otimes H_{[h(u) = a]}.
        \end{equation*}
    \item(Strong self-consistency): \label{item:self-improvement-self}
    	\begin{equation*}
		\sum_{h} \bra{\psi} H_h \ot H_h \ket{\psi} \geq \bra{\psi} H \ot I \ket{\psi} - \zeta.
	\end{equation*}
    \item(Boundedness): \label{item:self-improvement-boundedness} There exists a positive-semidefinite matrix~$Z$ such that
    \begin{equation*}
        \bra{\psi} Z \otimes I \ket{\psi} -\E_{\bu} \sum_a \bra{\psi}  A^{\bu}_{a} \otimes H_{[h(\bu)=a]} \ket{\psi} \leq \zeta
    \end{equation*}
    and for each $h \in \polyfunc{m}{q}{d}$,
	\begin{equation*}
	Z \geq \left(\E_{\bu} A^{\bu}_{h(\bu)}\right).
	\end{equation*}
\end{enumerate}
\end{lemma}

\subsection{A semidefinite program}

A key element in the proof of \Cref{lem:self-improvement-helper} will be a pair of primal and dual semidefinite programs. 
To define them, it will be convenient to introduce the notational shorthand
\begin{equation*}
A_g = \E_{\bu \sim \F_q^m} A^{\bu}_{g(\bu)}.
\end{equation*}
Then the primal is
\begin{align}
 \sup &\quad \sum_g \,\Tr(T_g \cdot A_g) \label{eq:primal-objective}\\
  \text{s.t.} &\quad   T_g \geq 0\qquad\forall g\in\polyfunc{m}{q}{d}\;,\notag\\
	&\quad \sum_g T_g \leq I,\notag
\end{align}
and the dual is
\begin{align}
  \inf &\quad \Tr(Z) \label{eq:dual-objective}\\
 \text{ s.t.} &\quad Z \geq A_g. \label{eq:dual-constraint}
%	&\quad Z \geq 0\;.\notag
\end{align}
We will prove that these two program are indeed dual to each other in \Cref{lem:sdp} below.

\begin{lemma}\label{lem:sdp}
The semidefinite programs~\eqref{eq:primal-objective} and~\eqref{eq:dual-objective} are dual to each other. 
Moreover there is an optimal pair of solutions $\{T_g\}$ to~\eqref{eq:primal-objective} and $Z$ to~\eqref{eq:dual-objective} such that $\sum_g T_g = I$ and  
\begin{equation}\label{eq:slater}
T_g Z  \,=\, T_g {A_g},\qquad\forall g\in \polyfunc{m}{q}{d}.
\end{equation}
\end{lemma}

\begin{proof}
To show that~\eqref{eq:primal-objective} and~\eqref{eq:dual-objective} are dual to each other,
we begin by rewriting the primal~\eqref{eq:primal-objective} in canonical form.
Let $r$ be the dimension of the space on which~$A$ acts.
Let $M = |\polyfunc{m}{q}{d}|$ be the number of polynomials with individual degree $d$. 
We will assume some ordering of these polynomials $g_1, \ldots, g_M \in \polyfunc{m}{q}{d}$ which is allowed to be arbitrary.
Consider the following semidefinite program:
\begin{align}
\sup &\quad \Tr(C^\dagger X) \label{eq:primal-canonical}\\
  \text{s.t.} &\quad   \Tr(D_{ij}^\dagger X) = b_{ij} \qquad\forall i,j\in\{1,\ldots,r\},\nonumber\\
	&\quad X \geq 0,\nonumber
\end{align}
where the variables~$C$, $D_{ij}$, and $b_{ij}$ are defined as follows:
\begin{equation*}
C = \sum_{i = 1}^M \ket{i}\bra{i} \otimes A_{g_i}
\end{equation*}
\begin{equation*}
\forall i,j \in [r], \quad
D_{ij} = \sum_{k=1}^{M+1} \ket{k}\bra{k} \otimes \ket{i}\bra{j},
\qquad
b_{ij}  = \left\{\begin{array}{rl}
		1 & \text{if } i = j,\\
		0 & \text{otherwise}.
		\end{array}\right.
\end{equation*}
We claim that \eqref{eq:primal-canonical} is equivalent to \eqref{eq:primal-objective}.
To see this, let
\begin{equation*}
X = \sum_{i, j =1}^{M+1} \ket{i}\bra{j} \otimes X_{ij}
\end{equation*}
be a feasible solution to \eqref{eq:primal-canonical}.
Then because $X \geq 0$, $X_{i i} \geq 0$ for each $i \in [M+1]$. In addition
\begin{align*}
\Tr(D_{i_1 i_2}^\dagger X)
&= \sum_k \sum_{j_1 j_2} \Tr((\ket{k}\bra{k} \otimes \ket{i_1}\bra{i_2})\cdot(\ket{j_1} \bra{j_2} \otimes X_{j_1 j_2}))\\
&= \sum_k  \bra{i_2} X_{k k} \ket{i_1}\\
&= \left\{\begin{array}{rl}
		1 & \text{if } i_1 = i_2,\\
		0 & \text{otherwise}.
		\end{array}\right.
\end{align*}
This is equivalent to the statement $\sum_{i=1}^{M+1} X_{i i} = I$, which implies that $\sum_{i=1}^M X_{i i} \leq I$.
Finally, the objective value \eqref{eq:primal-canonical} is
\begin{align*}
\Tr(C^\dagger X)
& = \sum_{i = 1}^M  \sum_{j,k = 1}^{M+1} \Tr((\ket{i}\bra{i} \otimes A_{g_i}^\dagger) \cdot (\ket{j}\bra{k} \otimes X_{j, k}))\\
&=  \sum_{i = 1}^M \Tr(A_{g_i}^\dagger \cdot  X_{i, i})\\
&=  \sum_{i = 1}^M \Tr(A_{g_i} \cdot  X_{i, i}). \tag{because~$A$ is Hermitian}
\end{align*}
As a result, setting $T_{g_i} = X_{i i}$ for each $i \in [M]$ gives a feasible solution to the original semidefinite program~\eqref{eq:primal-objective}
with a matching objective value. A similar transformation allows us to convert solutions of~\eqref{eq:primal-objective} to~\eqref{eq:primal-canonical}.

The dual of \eqref{eq:primal-canonical} is 
\begin{align}
\inf &\quad \sum_i\, z_{ij} b_{ij} \label{eq:dual-canonical}\\
  \text{s.t.} &\quad  \sum_{i,j} \,z_{ij} D_{ij} \geq C.\label{eq:dual-canonical-constraint}
\end{align}
We claim that \eqref{eq:dual-canonical} is equivalent to \eqref{eq:dual-objective}.
To see this, we first calculate
\begin{align*}
\sum_{i,j=1}^r \,z_{ij} D_{ij}
= \sum_{i,j=1}^r z_{ij} \cdot \sum_{k=1}^{M+1} \ket{k}\bra{k} \otimes \ket{i}\bra{j}
&= \sum_{k=1}^{M+1} \ket{k}\bra{k} \otimes \left(\sum_{i, j=1}^r z_{ij}\ket{i}\bra{j}\right)\\
&=: \sum_{k=1}^{M+1}\ket{k}\bra{k} \otimes Z.
\end{align*}
Then the constraint \eqref{eq:dual-canonical-constraint} states that
\begin{equation*}
\sum_{i=1}^{M+1}\ket{i}\bra{i} \otimes Z \geq C = \sum_{i = 1}^M \ket{i}\bra{i} \otimes A_{g_i},
\end{equation*}
which is equivalent to the statement that $Z \geq A_{g_i}$ for all $i \in [M]$ and $Z \geq 0$.
In other words, $Z$ is a feasible solution to the original dual SDP \eqref{eq:dual-objective}, with value
\begin{equation*}
\Tr(Z) = \sum_{i=1}^r Z_{i i} = \sum_{i,j=1}^r b_{i j} Z_{i j} =  \sum_{i,j=1}^r b_{i j} z_{i j},
\end{equation*}
the same value as in \eqref{eq:dual-canonical}.
Hence, the two dual programs are the same as well,
which implies that \eqref{eq:primal-objective} and \eqref{eq:dual-objective} form a primal/dual pair.

To show that a primal/dual pair satisfies \emph{strong duality}, i.e.\ that their optimum values are the same,
we use Slater's condition~\cite[Section~5.2.3]{cvxbook} and show that they satisfy \emph{strict feasibility}, which means both have a feasible solution which is positive \emph{definite}
that satisfies all constraints with a strict inequality.
It can be checked that the following two solutions to \eqref{eq:primal-objective} and \eqref{eq:dual-objective} satisfy this property:
\begin{equation*}
\forall g \in \polyfunc{m}{q}{d}, \quad T_g = \frac{1}{2M} \cdot I,
\qquad\qquad\qquad
Z = 2I.
\end{equation*}
Because they satisfy strong duality, their optimal solutions satisfy the \emph{complementary slackness} condition (see~\cite{AHO97}).
If $X$ and $(z_{ij})$ are an optimal pair of solutions to~\eqref{eq:primal-canonical} and~\eqref{eq:dual-canonical} respectively,  this implies that
\begin{equation}\label{eq:complementary-slackness}
X\Big( \sum_{i,j} \,z_{ij} D_{ij} - C \Big) \,=\, 0 \;.
\end{equation}
Clearly for any optimal pair $(X,z_{ij})$ we can assume without loss of generality that $X$ is block-diagonal.
Then if we translate \eqref{eq:complementary-slackness} back to the variables $\{T_g\}$ and $Z$, we get
\begin{align*}
0 = X\Big( \sum_{i,j} \,z_{ij} D_{ij} - C \Big)
&= \sum_{i=1}^{M+1} \ket{i}\bra{i} \otimes X_{ii} \cdot \bigg(\sum_{i=1}^{M+1} \ket{i}\bra{i} \otimes Z - \sum_{i=1}^M \ket{i}\bra{i}\otimes A_{g_i}\bigg)\\
&= \sum_{i=1}^{M} \ket{i}\bra{i} \otimes (X_{ii} \cdot (Z - A_{g_i})) + \ket{M+1}\bra{M+1} \otimes (X_{M+1,M+1} \cdot Z).
\end{align*}
This implies that $X_{ii} \cdot (Z - A_{g_i}) = 0$ for $i \in [M]$ and $X_{M+1, M+1} = 0$.
Translating back to the variables $\{T_g\}$ and $Z$, this gives $T_g(Z-A_g)=0$ for all $g\in\polyfunc{m}{q}{d}$ and $\sum_g T_g = I$. 
\end{proof}


\subsection{Proof of \Cref{lem:self-improvement-helper}}
We let $T = \{T_g\}$ and $Z$ be the optimal solutions to the SDPs~\eqref{eq:primal-objective} and \eqref{eq:dual-objective} respectively
given by \Cref{lem:sdp}.
Then~$T$ is a measurement, and
\begin{align}
 \forall g \in \polyfunc{m}{q}{d}: \qquad
 Z &\geq (\E_{\bu} A^{\bu}_{g(\bu)}),  \label{eq:Z-greater-than-A}\\
T_g \cdot Z &= T_g \cdot (\E_{\bu} A^{\bu}_{g(\bu)}).\label{eq:swap-Z-for-A}
\end{align}
For each $u \in \F_q^m$, define $H^u = \{H^u_h\}_{h \in \polyfunc{m}{q}{d}}$ as
\begin{equation*}
    H^u_h := A^u_{h(u)}\cdot T_h \cdot A^u_{h(u)}.
\end{equation*}
Let $u \in \F_q^m$. Then
\begin{align*}
    \sum_{h \in \polyfunc{m}{q}{d}} H^u_h
    & = \sum_{h \in \polyfunc{m}{q}{d}} A^u_{h(u)} \cdot T_h \cdot A^u_{h(u)}\\
    & = \sum_{a \in \F_q}A^u_{a} \cdot \bigg(\sum_{h:h(u) = a}  T_h\bigg) \cdot A^u_{a}\\
    & \leq \sum_{a \in \F_q} (A^u_{a})^2\tag{because~$T$ is a measurement}\\
    & = \sum_{a \in \F_q} A^u_{a}\tag{$A$ is projective}\\
    & = I.
\end{align*}
Hence, $H^u$ is a sub-measurement,
and therefore $H^u \in \polysub{m}{q}{d}$.
Next, define $H = \{H_h\}_{h \in \polyfunc{m}{q}{d}}$ as
\begin{equation*}
H_h := \E_{\bu \sim \F_q^m}H^{\bu}_h.
\end{equation*}
Then
\begin{equation*}
\sum_{h \in \polyfunc{m}{q}{d}} H_h
= \E_{\bu \sim \F_q^m} \sum_{h \in \polyfunc{m}{q}{d}} H^{\bu}_h
\leq I. \tag{$H^u$ is a sub-measurement}
\end{equation*}
Hence, $H$ is a sub-measurement,
and therefore $H \in \polysub{m}{q}{d}$.

Set
\begin{equation*}
\zeta_{\mathrm{variance}} = 24m\cdot\Big(\eps + \delta + \frac{md}{q}\Big)
\end{equation*}
to be the error in \Cref{eq:local-variance-of-points-equation}.
Prior to showing that~$H$ satisfies \Cref{item:self-improvement-A-consistency,item:self-improvement-completeness,item:self-improvement-self,item:self-improvement-boundedness}, we will prove the following technical lemma.

\begin{lemma}\label{lem:add-in-u}
Suppose $M = \{M^u_o\}$ is a sub-measurement with outcomes in some set $\calO$.
For each $u \in \F_q^m$, let $S_u$ be a subset of $\calO \otimes \polyfunc{m}{q}{d}$.
Then
\begin{equation*}
\E_{\bu \sim \F_q^m} \sum_{(o, h) \in S_{\bu}} \bra{\psi} M^{\bu}_o \otimes H_h \ket{\psi}
\approx_{4\sqrt{\zeta_{\mathrm{variance}}}} \E_{\bu \sim \F_q^m} \sum_{(o, h) \in S_{\bu}} \bra{\psi} (A^{\bu}_{h(\bu)} \cdot M^{\bu}_o \cdot A^{\bu}_{h(\bu)}) \otimes T_h \ket{\psi}.
\end{equation*}
\end{lemma}
\begin{proof}
We begin by expanding
\begin{align}
\E_{\bu \sim \F_q^m} \sum_{(o, h) \in S_{\bu}} \bra{\psi} M^{\bu}_o \otimes H_h \ket{\psi}
&= \E_{\bu, \bv} \sum_{(o, h) \in S_{\bu}} \bra{\psi} M^{\bu}_o \otimes H^{\bv}_h \ket{\psi}\nonumber\\
&= \E_{\bu, \bv} \sum_{(o, h) \in S_{\bu}} \bra{\psi} M^{\bu}_o \otimes (A^{\bv}_{h(\bv)} \cdot T_h \cdot A^{\bv}_{h(\bv)}) \ket{\psi}.\label{eq:expand-that-H}
\end{align}
We claim that
\begin{equation}\label{eq:move-one}
\eqref{eq:expand-that-H}
\approx_{\sqrt{2\delta}}
\E_{\bu, \bv} \sum_{(o, h) \in S_{\bu}} \bra{\psi} (A^{\bv}_{h(\bv)} \cdot M^{\bu}_o) \otimes (T_h \cdot A^{\bv}_{h(\bv)}) \ket{\psi}.
\end{equation}
To show this, we bound the magnitude of the difference.
\begin{multline}
\Big|\E_{\bu, \bv} \sum_{(o, h) \in S_{\bu}} \bra{\psi} (A^{\bv}_{h(\bv)} \otimes I - I \otimes A^{\bv}_{h(\bv)}) \cdot (M^{\bu}_o \otimes (T_h \cdot A^{\bv}_{h(\bv)}))\ket{\psi}\Big|\\
\leq
\sqrt{\E_{\bu, \bv} \sum_{(o, h) \in S_{\bu}} \bra{\psi} (A^{\bv}_{h(\bv)} \otimes I - I \otimes A^{\bv}_{h(\bv)}) \cdot (M^{\bu}_o \otimes T_h) \cdot (A^{\bv}_{h(\bv)} \otimes I - I \otimes A^{\bv}_{h(\bv)}) \ket{\psi}}\\
\cdot \sqrt{\E_{\bu, \bv} \sum_{(o, h) \in S_{\bu}} \bra{\psi}  M^{\bu}_o \otimes (A^{\bv}_{h(\bv)} \cdot T_h\cdot  A^{\bv}_{h(\bv)}) \ket{\psi}}.\label{eq:move-one-cauchy-schwarz}
\end{multline}
The term inside the first square root is
\begin{align*}
&\E_{\bu, \bv} \sum_{a \in \F_q} \bra{\psi} (A^{\bv}_a \otimes I - I \otimes A^{\bv}_a) \cdot \bigg( \sum_{(o, h) \in S_{\bu},h(\bv) = a}M^{\bu}_o \otimes T_h\bigg) \cdot (A^{\bv}_a \otimes I - I \otimes A^{\bv}_a) \ket{\psi}\\
\leq~&\E_{\bu, \bv} \sum_{a \in \F_q} \bra{\psi} (A^{\bv}_a \otimes I - I \otimes A^{\bv}_a)^2 \ket{\psi}, \tag{because $M^{\bu}$ and $T_h$ are sub-measurements}
\end{align*}
which is at most~$2\delta$ by \Cref{prop:simeq-to-approx} and because~$A$ is $\delta$-self-consistent.
The term inside the second square root is
\begin{equation*}
\E_{\bu, \bv} \sum_{(o, h) \in S_{\bu}} \bra{\psi}  M^{\bu}_o \otimes H^{\bv}_h \ket{\psi},
\end{equation*}
which is at most~$1$ because $M^{\bu}$ and $H^{\bv}$ are sub-measurements.
Next, we claim that
\begin{equation}\label{eq:move-another}
\eqref{eq:move-one}
\approx_{\sqrt{2\delta}}
\E_{\bu, \bv} \sum_{(o, h) \in S_{\bu}} \bra{\psi} (A^{\bv}_{h(\bv)} \cdot M^{\bu}_o \cdot A^{\bv}_{h(\bv)}) \otimes T_h  \ket{\psi}.
\end{equation}
To show this, we bound the magnitude of the difference.
\begin{align}\label{eq:move-another-cauchy-schwarz}
&\Big|\E_{\bu, \bv} \sum_{(o, h) \in S_{\bu}} \bra{\psi} ((A^{\bv}_{h(\bv)} \cdot M^{\bu}_o) \otimes T_h) \cdot (A^{\bv}_{h(\bv)} \otimes I - I \otimes A^{\bv}_{h(\bv)})  \ket{\psi}\Big|\nonumber\\
&\leq\sqrt{\E_{\bu, \bv} \sum_{(o, h) \in S_{\bu}} \bra{\psi} (A^{\bv}_{h(\bv)} \cdot M^{\bu}_o \cdot A^{\bv}_{h(\bv)}) \otimes T_h  \ket{\psi}}\nonumber\\
&\quad\cdot \sqrt{\E_{\bu, \bv} \sum_{(o, h) \in S_{\bu}} \bra{\psi} (A^{\bv}_{h(\bv)} \otimes I - I \otimes A^{\bv}_{h(\bv)}) \cdot (M^{\bu}_o \otimes T_h) \cdot (A^{\bv}_{h(\bv)} \otimes I - I \otimes A^{\bv}_{h(\bv)}) \ket{\psi}}.
\end{align}
The term inside the first square root is
\begin{align*}
&\E_{\bu, \bv} \sum_h \bra{\psi} (A^{\bv}_{h(\bv)} \cdot \bigg(\sum_{o:(o, h) \in S_{\bu}} M^{\bu}_o \bigg)\cdot A^{\bv}_{h(\bv)}) \otimes T_h  \ket{\psi}\\
\leq~&\E_{\bu, \bv} \sum_h \bra{\psi} (A^{\bv}_{h(\bv)})^2 \otimes T_h  \ket{\psi}\tag{because~$M$ is a sub-measurement}\\
\leq~&\E_{\bu, \bv} \sum_h \bra{\psi} I \otimes T_h  \ket{\psi}\tag{because~$A^{\bv}_{h(\bv)} \leq I$}\\
=~&1.\tag{because~$T$ is a measurement}
\end{align*}
As for the term inside the second square root, it is equal to the term inside the first square root in \Cref{eq:move-one-cauchy-schwarz}, which we showed was at most $2\delta$.

Having moved both $A$'s to the left-hand side, we want to show that
\begin{equation}\label{eq:change-one}
\eqref{eq:move-another}
\approx_{\sqrt{\zeta_{\mathrm{variance}}}}
\E_{\bu, \bv} \sum_{(o, h) \in S_{\bu}} \bra{\psi} (A^{\bu}_{h(\bu)} \cdot M^{\bu}_o \cdot A^{\bv}_{h(\bv)}) \otimes T_h  \ket{\psi}.
\end{equation} 
To show this, we bound the magnitude of the difference.
\begin{multline}\label{eq:change-one-cauchy-schwarz}
\Big|\E_{\bu, \bv} \sum_{(o, h) \in S_{\bu}} \bra{\psi} ((A^{\bv}_{h(\bv)} - A^{\bu}_{h(\bu)}) \cdot M^{\bu}_o \cdot A^{\bv}_{h(\bv)}) \otimes T_h  \ket{\psi}\Big|\\
\leq
\sqrt{\E_{\bu, \bv} \sum_{(o, h) \in S_{\bu}} \bra{\psi} ((A^{\bv}_{h(\bv)} - A^{\bu}_{h(\bu)}) \cdot M^{\bu}_o \cdot (A^{\bv}_{h(\bv)} - A^{\bu}_{h(\bu)})) \otimes T_h  \ket{\psi}}\\
\cdot \sqrt{\E_{\bu, \bv} \sum_{(o, h) \in S_{\bu}} \bra{\psi} (A^{\bv}_{h(\bv)} \cdot M^{\bu}_o \cdot A^{\bv}_{h(\bv)}) \otimes T_h  \ket{\psi}}.
\end{multline}
The term inside the first square root is
\begin{align*}
&\E_{\bu, \bv} \sum_h \bra{\psi} ((A^{\bv}_{h(\bv)} - A^{\bu}_{h(\bu)}) \cdot \bigg( \sum_{o:(o, h) \in S_{\bu}}M^{\bu}_o\bigg) \cdot (A^{\bv}_{h(\bv)} - A^{\bu}_{h(\bu)})) \otimes T_h  \ket{\psi}\\
\leq~&\E_{\bu, \bv} \sum_h \bra{\psi} (A^{\bv}_{h(\bv)} - A^{\bu}_{h(\bu)})^2 \otimes T_h  \ket{\psi}. \tag{because~$M^{\bu}$ is a sub-measurement}
\end{align*}
But $T \in \polysub{m}{q}{d}$, and so by \Cref{lem:global-variance-of-points} this expression is at most $\zeta_{\mathrm{variance}}$.
As for the term inside the second square root, it is equal to the term inside the first square root in \Cref{eq:move-another-cauchy-schwarz}, which we showed was at most~$1$.
Finally, we want to show that
\begin{equation}\label{eq:change-another}
\eqref{eq:change-one}
\approx_{\sqrt{\zeta_{\mathrm{variance}}}}
\E_{\bu, \bv} \sum_{(o, h) \in S_{\bu}} \bra{\psi} (A^{\bu}_{h(\bu)} \cdot M^{\bu}_o \cdot A^{\bu}_{h(\bu)}) \otimes T_h  \ket{\psi}.
\end{equation} 
To show this, we bound the magnitude of the difference.
\begin{multline*}
\Big|
\E_{\bu, \bv} \sum_{(o, h) \in S_{\bu}} \bra{\psi} (A^{\bu}_{h(\bu)} \cdot M^{\bu}_o \cdot (A^{\bv}_{h(\bv)} - A^{\bu}_{h(\bu)})) \otimes T_h  \ket{\psi}\Big|\\
\leq
\sqrt{\E_{\bu, \bv} \sum_{(o, h) \in S_{\bu}} \bra{\psi} (A^{\bu}_{h(\bu)} \cdot M^{\bu}_o \cdot A^{\bu}_{h(\bu)}) \otimes T_h  \ket{\psi}}\\
\cdot \sqrt{\E_{\bu, \bv} \sum_{(o, h) \in S_{\bu}} \bra{\psi} ((A^{\bv}_{h(\bv)} - A^{\bu}_{h(\bu)}) \cdot M^{\bu}_o \cdot (A^{\bv}_{h(\bv)} - A^{\bu}_{h(\bu)})) \otimes T_h  \ket{\psi}}.
\end{multline*}
The term inside the first square root is
\begin{align*}
&\E_{\bu, \bv} \sum_h \bra{\psi} (A^{\bu}_{h(\bu)} \cdot \bigg(\sum_{o:(o, h) \in S_{\bu}} M^{\bu}_o \bigg)\cdot A^{\bu}_{h(\bu)}) \otimes T_h  \ket{\psi}\\
\leq~&\E_{\bu, \bv} \sum_h \bra{\psi} (A^{\bu}_{h(\bu)})^2 \otimes T_h  \ket{\psi}\tag{because~$M$ is a sub-measurement}\\
\leq~&\E_{\bu, \bv} \sum_h \bra{\psi} I \otimes T_h  \ket{\psi}\tag{because~$A^{\bu}_{h(\bu)} \leq I$}\\
=~&1.\tag{because~$T$ is a measurement}
\end{align*}
As for the term inside the second square root, it is equal to the term inside the first square root in \Cref{eq:change-one-cauchy-schwarz}, which we showed was at most $\zeta_{\mathrm{variance}}$.
This concludes the proof with an error of $2\sqrt{2\delta} + 2\sqrt{\zeta_{\mathrm{variance}}}$.
The lemma now follows by observing that $2\delta \leq \zeta_{\mathrm{variance}}$.
\end{proof}

We now show that~$H$ satisfies \Cref{item:self-improvement-completeness,item:self-improvement-A-consistency,item:self-improvement-self,item:self-improvement-boundedness}.



\vspace{\baselineskip}
\noindent
\emph{Proof of \Cref{item:self-improvement-completeness} (Completeness).}
The completeness of~$H$ is
\begin{align}
\sum_h \bra{\psi} H_h \otimes I \ket{\psi}
& = \E_{\bu} \sum_h \bra{\psi} H^{\bu}_h \otimes I \ket{\psi}\nonumber\\
&= \E_{\bu} \sum_h \bra{\psi} (A^{\bu}_{h(\bu)} \cdot T_h \cdot A^{\bu}_{h(\bu)}) \otimes I \ket{\psi}\nonumber\\
&= \E_{\bu} \sum_a \bra{\psi} (A^{\bu}_{a} \cdot T_{[h(\bu) = a]} \cdot A^{\bu}_{a}) \otimes I \ket{\psi}.\label{eq:bracketize-the-expression}
\end{align}
We claim that
\begin{equation}\label{eq:yet-another-move-a}
\eqref{eq:bracketize-the-expression}
\approx_{2\sqrt{\delta}}
\E_{\bu} \sum_a \bra{\psi} ( T_{[h(\bu)=a]} \cdot A^{\bu}_{a}) \otimes A^{\bu}_{a} \ket{\psi}.
\end{equation}
To show this, we bound the magnitude of the difference
\begin{multline*}
\Big|\E_{\bu} \sum_a \bra{\psi} (A^{\bu}_{a} \otimes I - I \otimes A^{\bu}_{a}) \cdot (( T_{[h(\bu) = a]} \cdot A^{\bu}_{a}) \otimes I) \ket{\psi}\Big|\\
\leq
\sqrt{\E_{\bu} \sum_a \bra{\psi} (A^{\bu}_{a} \otimes I - I \otimes A^{\bu}_{a})^2 \ket{\psi}}
\cdot \sqrt{\E_{\bu} \sum_a \bra{\psi} (A^{\bu}_{a} \cdot  T^2_{[h(\bu) = a]} \cdot A^{\bu}_{a}) \otimes I \ket{\psi}}.
\end{multline*}
The expression inside the first square root is at most $2\delta \leq 4 \delta$
by \Cref{prop:simeq-to-approx} and the self-consistency of~$A$,
and the expression inside the second square root is at most~$1$ because $T_{[h(\bu) = a]} \leq I$.
Next, we claim that 
\begin{equation}\label{eq:mysterious-case-of-the-disappearing-a}
\eqref{eq:yet-another-move-a}
\approx_{\sqrt{\delta}}
\E_{\bu} \sum_a \bra{\psi} ( T_{[h(\bu)=a]} \cdot A^{\bu}_{a}) \otimes I \ket{\psi}.
\end{equation}
To show this, we bound the magnitude of the difference.
\begin{multline*}
\Big|\E_{\bu} \sum_a \bra{\psi} ( T_{[h(\bu)=a]} \cdot A^{\bu}_{a}) \otimes (I-A^{\bu}_{a}) \ket{\psi}\Big|\\
\leq
\sqrt{\E_{\bu} \sum_a \bra{\psi} ( T_{[h(\bu)=a]})^2 \otimes I \ket{\psi}}
\cdot \sqrt{\E_{\bu} \sum_a \bra{\psi} ( A^{\bu}_{a})^2 \otimes (I-A^{\bu}_{a})^2 \ket{\psi}}.
\end{multline*}
The expression inside the first square root is at most~$1$ because~$T$ is a sub-measurement.
By the projectivity of~$A$, the expression inside the second square root is equal to
\begin{equation*}
\E_{\bu} \sum_a \bra{\psi} ( A^{\bu}_{a}) \otimes (I-A^{\bu}_{a}) \ket{\psi}
= \E_{\bu} \sum_{a\neq b} \bra{\psi} A^{\bu}_{a} \otimes A^{\bu}_{b} \ket{\psi},
\end{equation*}
which is at most~$\delta$ by the self-consistency of~$A$.
We can rewrite \Cref{eq:mysterious-case-of-the-disappearing-a} as
\begin{align*}
&\E_{\bu} \sum_h \bra{\psi} ( T_{h} \cdot A^{\bu}_{h(\bu)}) \otimes I \ket{\psi}\\
 =~&  \sum_h \bra{\psi} ( T_{h} \cdot\E_{\bu} A^{\bu}_{h(\bu)}) \otimes I \ket{\psi}\\
 =~&  \sum_h \bra{\psi} ( T_{h} \cdot Z) \otimes I \ket{\psi}\tag{by \Cref{eq:swap-Z-for-A}}\\
=~&  \bra{\psi}  Z \otimes I \ket{\psi}.\tag{because~$T$ is a measurement}
\end{align*}
We pause and record what we have shown so far:
\begin{equation}\label{eq:gonna-use-this-later-H-versus-Z}
\sum_h \bra{\psi} H_h \otimes I \ket{\psi} \geq \bra{\psi}  Z \otimes I \ket{\psi} - 3\sqrt{\delta}.
\end{equation}
At this point, we can lower-bound 
\begin{align*}
\bra{\psi}  Z \otimes I \ket{\psi}
\geq~& \sum_g \bra{\psi} Z \otimes G_g \ket{\psi} \tag{because~$G$ is a sub-measurement}\\
\geq~& \sum_g \bra{\psi} (\E_{\bu} A^{\bu}_{g(\bu)}) \otimes G_g \ket{\psi} \tag{by \Cref{eq:Z-greater-than-A}}\\
=~& \E_{\bu} \sum_a \bra{\psi}  A^{\bu}_{a} \otimes G_{[g(\bu)=a]} \ket{\psi}\\
\geq~& 1-\nu.\tag{by \Cref{item:self-improvement-G-consistency} and \Cref{prop:simeq-for-measurements}}
\end{align*}
Thus, the completeness is at least $1  - \nu - 3\sqrt{\delta}$. 
The proof follows from noting that $3\sqrt{\delta} \leq \zeta$.

\vspace{\baselineskip}
\noindent
\emph{Proof of \Cref{item:self-improvement-A-consistency} (Consistency with~$A$).}
The inconsistency of~$H$ with~$A$ is
\begin{equation}\label{eq:consistency-with-A-baby-step}
\E_{\bu \sim \F_q^m} \sum_{a \neq b} \bra{\psi} A^{\bu}_a \otimes H_{[h(\bu) = b]} \ket{\psi}
= \E_{\bu \sim \F_q^m} \sum_{a, h: h(\bu) \neq a} \bra{\psi} A^{\bu}_a \otimes H_h \ket{\psi}.
\end{equation}
We now apply \Cref{lem:add-in-u} to the right-hand side of \Cref{eq:consistency-with-A-baby-step}.
To do so, we set $\calO = \F_q^m$, $M = A$, and $S_u = \{(a, h) : h(u) \neq a\}$.
Then \Cref{lem:add-in-u} implies that
\begin{equation*}
\eqref{eq:consistency-with-A-baby-step}
\approx_{4 \sqrt{\zeta_{\mathrm{variance}}}}
\E_{\bu \sim \F_q^m} \sum_{a, h: h(\bu) \neq a} \bra{\psi} (A^{\bu}_{h(\bu)} \cdot A^{\bu}_a  \cdot A^{\bu}_{h(\bu)})\otimes T_h \ket{\psi},
\end{equation*}
which is equal to~$0$ because~$A$ is projective.
This implies that
\begin{equation}\label{eq:explicit-bound-for-A-consistency}
\E_{\bu \sim \F_q^m} \sum_{a \neq b} \bra{\psi} A^{\bu}_a \otimes H_{[h(\bu) = b]} \ket{\psi} \leq 4 \sqrt{\zeta_{\mathrm{variance}}}.
\end{equation}
The proof follows from noting that
\begin{align*}
4\sqrt{\zeta_{\mathrm{variance}}}
&= 4 \cdot \sqrt{24 m \cdot\Big(\eps + \delta + \frac{md}{q}\Big)}\\
&\leq 20 m \cdot \Big(\eps^{1/2} + \delta^{1/2} + (d/q)^{1/2}\Big)\\
&\leq \zeta.
\end{align*}

\vspace{\baselineskip}
\noindent
\emph{Proof of \Cref{item:self-improvement-self} (Strong self-consistency).}
We begin by recording the following facts which follow from the definition of $H^u_h$ and the projectivity of~$A$:
\begin{align}
H^u_h &= A^u_{h(u)} \cdot T_h \cdot A^u_{h(u)} = A^u_{h(u)} \cdot H^u_h \cdot A^u_{h(u)},\label{eq:h-sandwich}\\
A^u_{h(u)} \cdot H^u_{h'} \cdot A^u_{h(u)} &=  H^u_{h'} \cdot A^u_{h(u)} = (A^u_{h(u)} \cdot T_{h'} \cdot A^u_{h(u)}) \cdot \bone[h(u) = h'(u)].\label{eq:h-blt}
\end{align}
The strong self-consistency of~$H$ is
\begin{equation}\label{eq:self-consistency-baby-step}
\sum_{h\in\polyfunc{m}{q}{d}} \bra{\psi} H_h \otimes H_h \ket{\psi}
= \E_{\bu \sim \F_q^m} \sum_{h \in \polyfunc{m}{q}{d}} \bra{\psi} H_h^{\bu} \otimes H_h \ket{\psi}.
\end{equation}
We now apply \Cref{lem:add-in-u} to the right-hand side of \Cref{eq:self-consistency-baby-step}.
To do so, we set
\begin{equation*}
\text{$\calO =\polyfunc{m}{q}{d}$, $M = H$, and $S_u = \{(h, h) : h \in \polyfunc{m}{q}{d}\}$.}
\end{equation*}
Then \Cref{lem:add-in-u} implies that
\begin{equation}\label{eq:release-the-kraken}
\eqref{eq:self-consistency-baby-step}
\approx_{4 \sqrt{\zeta_{\mathrm{variance}}}}
\E_{\bu \sim \F_q^m} \sum_{h \in \polyfunc{m}{q}{d}} \bra{\psi} (A^{\bu}_{h(\bu)} \cdot H_h^{\bu} \cdot A^{\bu}_{h(\bu)}) \otimes T_h \ket{\psi}.
\end{equation}
Having placed two $A$'s on the left-hand side, we want to show that
\begin{equation}\label{eq:threw-in-h-prime}
\eqref{eq:release-the-kraken}
\approx_{2\sqrt{\zeta_{\mathrm{variance}}} + \frac{md}{q}}
\E_{\bu \sim \F_q^m} \sum_{h, h' \in \polyfunc{m}{q}{d}} \bra{\psi} (A^{\bu}_{h(\bu)} \cdot H^{\bu}_{h'} \cdot A^{\bu}_{h(\bu)}) \otimes T_h \ket{\psi}.
\end{equation}
We note that \eqref{eq:threw-in-h-prime} is at least as big as \eqref{eq:release-the-kraken}. Thus, we want to upper-bound
\begin{align}
\eqref{eq:threw-in-h-prime}-\eqref{eq:release-the-kraken}
& = \E_{\bu \sim \F_q^m} \sum_{h \neq h'} \bra{\psi}  (A^{\bu}_{h(\bu)} \cdot H^{\bu}_{h'} \cdot A^{\bu}_{h(\bu)}) \otimes T_h \ket{\psi}\nonumber\\
& = \E_{\bu \sim \F_q^m} \sum_{h \neq h'} \bra{\psi}  (A^{\bu}_{h(\bu)} \cdot T_{h'} \cdot A^{\bu}_{h(\bu)}) \otimes T_h \ket{\psi} \cdot \bone[h(\bu) = h'(\bu)],\label{eq:added-indicator}
\end{align}
where the second equality is by~\Cref{eq:h-blt}.
To do this, we first show that
\begin{equation}\label{eq:swapped-u-for-v}
\eqref{eq:added-indicator}
\approx_{\sqrt{\zeta_{\mathrm{variance}}}}
\E_{\bu,\bv} \sum_{h \neq h'} \bra{\psi}  (A^{\bv}_{h(\bv)} \cdot T_{h'} \cdot A^{\bu}_{h(\bu)}) \otimes T_h \ket{\psi} \cdot \bone[h(\bu) = h'(\bu)].
\end{equation}
To show this, we bound the magnitude of the difference.
\begin{multline}
\Big|\E_{\bu,\bv} \sum_{h \neq h'} \bra{\psi} ( (A^{\bu}_{h(\bu)} - A^{\bv}_{h(\bv)}) \cdot T_{h'} \cdot A^{\bu}_{h(\bu)}) \otimes T_h \ket{\psi} \cdot \bone[h(\bu) = h'(\bu)]\Big|\\
\leq
\sqrt{\E_{\bu,\bv} \sum_{h \neq h'} \bra{\psi}  ( (A^{\bu}_{h(\bu)} - A^{\bv}_{h(\bv)}) \cdot T_{h'} \cdot (A^{\bu}_{h(\bu)} - A^{\bv}_{h(\bv)})) \otimes T_h\ket{\psi}}\\
\cdot \sqrt{\E_{\bu,\bv} \sum_{h \neq h'} \bra{\psi} (A^{\bu}_{h(\bu)} \cdot T_{h'} \cdot A^{\bu}_{h(\bu)}) \otimes T_h\ket{\psi} \cdot \bone[h(\bu) = h'(\bu)]}. \label{eq:swapped-u-for-cauchy-schwarz}
\end{multline}
The term inside the first square root is
\begin{align*}
&\E_{\bu,\bv} \sum_{h} \bra{\psi} ( (A^{\bu}_{h(\bu)} - A^{\bv}_{h(\bv)}) \cdot \bigg(\sum_{h' \neq h} T_{h'}\bigg) \cdot (A^{\bu}_{h(\bu)} - A^{\bv}_{h(\bv)})) \otimes T_h\ket{\psi}\\
\leq~& \E_{\bu,\bv} \sum_{h} \bra{\psi}  (A^{\bu}_{h(\bu)} - A^{\bv}_{h(\bv)})^2 \otimes T_h\ket{\psi},\tag{because $T$ is a sub-measurement}
\end{align*}
But $T \in \polyfunc{m}{q}{d}$, and so by \Cref{lem:global-variance-of-points} this expression is at most~$\zeta_{\mathrm{variance}}$.
The term inside the second square root is equal to
\begin{equation*}
\E_{\bu,\bv} \sum_{h \neq h'} \bra{\psi}  H^{\bu}_{h'} \otimes T_h \ket{\psi} \cdot \bone[h(\bu) = h'(\bu)],
\end{equation*}
which is at most~$1$ because $T$ and $H^{\bu}$ are sub-measurements.
Next, we show that
\begin{equation}\label{eq:swapped-u-for-v-this-time-it's-personal}
\eqref{eq:swapped-u-for-v}
\approx_{\sqrt{\zeta_{\mathrm{variance}}}}
\E_{\bu,\bv} \sum_{h \neq h'} \bra{\psi}  (A^{\bv}_{h(\bv)} \cdot T_{h'} \cdot A^{\bv}_{h(\bv)}) \otimes T_h \ket{\psi} \cdot \bone[h(\bu) = h'(\bu)].
\end{equation}
To show this, we bound the magnitude of the difference.
\begin{multline*}
\Big|\E_{\bu,\bv} \sum_{h \neq h'} \bra{\psi}  (A^{\bv}_{h(\bv)} \cdot T_{h'} \cdot (A^{\bu}_{h(\bu)} - A^{\bv}_{h(\bv)})) \otimes T_h \ket{\psi} \cdot \bone[h(\bu) = h'(\bu)]\Big|\\
\leq
\sqrt{\E_{\bu,\bv} \sum_{h \neq h'} \bra{\psi}  (A^{\bv}_{h(\bv)} \cdot T_{h'} \cdot A^{\bv}_{h(\bv)}) \otimes T_h\ket{\psi} \cdot \bone[h(\bu) = h'(\bu)]}\\
\cdot \sqrt{\E_{\bu,\bv} \sum_{h \neq h'} \bra{\psi} ( (A^{\bu}_{h(\bu)} - A^{\bv}_{h(\bv)}) \cdot T_{h'} \cdot (A^{\bu}_{h(\bu)} - A^{\bv}_{h(\bv)}))\otimes T_h\ket{\psi}}.
\end{multline*}
The term inside the first square root is at most
\begin{align}
\E_{\bu,\bv} \sum_{h \neq h'} \bra{\psi}  (A^{\bv}_{h(\bv)} \cdot T_{h'} \cdot A^{\bv}_{h(\bv)})\otimes T_h\ket{\psi}
& = \E_{\bu,\bv} \sum_{h} \bra{\psi}  (A^{\bv}_{h(\bv)} \cdot \bigg(\sum_{h' \neq h} T_{h'}\bigg) \cdot A^{\bv}_{h(\bv)}) \otimes T_h\ket{\psi}\\
& \leq \E_{\bu,\bv} \sum_{h} \bra{\psi}   (A^{\bv}_{h(\bv)})^2\otimes T_h\ket{\psi}\tag{$T$ is a sub-measurement}\\
& \leq \E_{\bu,\bv} \sum_{h} \bra{\psi}   I \otimes T_h\ket{\psi}\tag{because $A^{\bv}_{h(\bv)} \leq I$}\\
& \leq 1, \label{eq:gonna-use-this-later}
\end{align}
where the last step again uses the fact that~$T$ is a sub-measurement.
As for the term inside the second square root, it is equal to the term inside the first square root in \Cref{eq:swapped-u-for-cauchy-schwarz},
which we showed was at most $\zeta_{\mathrm{variance}}$.
Finally, \Cref{eq:swapped-u-for-v-this-time-it's-personal} is equal to
\begin{align*}
&\E_{\bv} \sum_{h \neq h'} \bra{\psi} (A^{\bv}_{h(\bv)} \cdot T_{h'} \cdot A^{\bv}_{h(\bv)}) \otimes T_h\ket{\psi} \cdot \E_{\bu}\bone[h(\bu) = h'(\bu)]\\
\leq~& \E_{\bv} \sum_{h \neq h'} \bra{\psi} (A^{\bv}_{h(\bv)} \cdot T_{h'} \cdot A^{\bv}_{h(\bv)}) \otimes T_h \ket{\psi} \cdot \frac{md}{q} \tag{by Schwartz-Zippel}\\
\leq~& \frac{md}{q}. \tag{by \Cref{eq:gonna-use-this-later}}
\end{align*}

By~\Cref{eq:h-blt}, \Cref{eq:threw-in-h-prime} is equal to
\begin{equation}\label{eq:delete-an-A}
\E_{\bu \sim \F_q^m} \sum_{h, h' \in \polyfunc{m}{q}{d}} \bra{\psi} (H^{\bu}_{h'} \cdot A^{\bu}_{h(\bu)}) \otimes T_h \ket{\psi}.
\end{equation}
Now, we show that
\begin{equation}\label{eq:swap-u-for-v-attack-of-the-clones}
\eqref{eq:delete-an-A}
\approx_{\sqrt{\zeta_{\mathrm{variance}}}}
\E_{\bu, \bv} \sum_{h, h' \in \polyfunc{m}{q}{d}} \bra{\psi} (H^{\bu}_{h'} \cdot A^{\bv}_{h(\bv)}) \otimes T_h \ket{\psi}.
\end{equation}
To show this, we bound the magnitude of the difference.
\begin{align*}
&\Big|\E_{\bu, \bv} \sum_{h, h' \in \polyfunc{m}{q}{d}} \bra{\psi} (H^{\bu}_{h'} \cdot (A^{\bu}_{h(\bu)} - A^{\bv}_{h(\bv)})) \otimes T_h \ket{\psi}\Big|\\
&\leq
\sqrt{\E_{\bu, \bv} \sum_{h, h' \in \polyfunc{m}{q}{d}} \bra{\psi} H^{\bu}_{h'} \otimes T_h\ket{\psi}}\\
&\quad\cdot \sqrt{\E_{\bu, \bv} \sum_{h, h' \in \polyfunc{m}{q}{d}} \bra{\psi} ((A^{\bu}_{h(\bu)} - A^{\bv}_{h(\bv)}) \cdot H^{\bu}_{h'} \cdot (A^{\bu}_{h(\bu)} - A^{\bv}_{h(\bv)})) \otimes T_h\ket{\psi}}.
\end{align*}
The expression inside the first square root is at most~$1$ because~$T$ and~$H^{\bu}$ are sub-measurements.
The term inside the second square root is
\begin{align*}
&\E_{\bu, \bv} \sum_{h} \bra{\psi} ((A^{\bu}_{h(\bu)} - A^{\bv}_{h(\bv)}) \cdot \bigg( \sum_{h'} H^{\bu}_{h'}\bigg) \cdot (A^{\bu}_{h(\bu)} - A^{\bv}_{h(\bv)})) \otimes T_h \ket{\psi}\\
\leq~& \E_{\bu, \bv} \sum_{h} \bra{\psi}  (A^{\bu}_{h(\bu)} - A^{\bv}_{h(\bv)})^2 \otimes T_h\ket{\psi}.
\tag{because $H^{\bu}$ is a sub-measurement}
\end{align*}
But $T \in \polysub{m}{q}{d}$, and so by \Cref{lem:global-variance-of-points} this expression is at most~$\zeta_{\mathrm{variance}}$.
Next, we show that
\begin{equation}\label{eq:move-over-v}
\eqref{eq:swap-u-for-v-attack-of-the-clones}
\approx_{\sqrt{2\delta}}
\E_{\bu, \bv} \sum_{h, h' \in \polyfunc{m}{q}{d}} \bra{\psi} H^{\bu}_{h'}  \otimes (T_h \cdot A^{\bv}_{h(\bv)}) \ket{\psi}.
\end{equation}
To show this, we bound the magnitude of the difference.
\begin{align*}
&\Big|\E_{\bu, \bv} \sum_{h, h' \in \polyfunc{m}{q}{d}} \bra{\psi}  (H^{\bu}_{h'}\otimes T_h) \cdot (A^{\bv}_{h(\bv)} \otimes I - I \otimes A^{\bv}_{h(\bv)}) \ket{\psi}\Big|\\
=~&\Big|\E_{\bu, \bv} \sum_{a \in \F_q}  \bra{\psi}  \bigg( \sum_{h'}H^{\bu}_{h'} \otimes \sum_{h:h(\bv) = a} T_h \bigg) \cdot (A^{\bv}_a \otimes I - I \otimes A^{\bv}_a) \ket{\psi}\Big|\\
\leq~&
\sqrt{\E_{\bu, \bv} \sum_{a \in \F_q}  \bra{\psi}  \bigg( \sum_{h'}H^{\bu}_{h'} \otimes \sum_{h:h(\bv) = a} T_h\bigg)^2\ket{\psi}}
\cdot \sqrt{\E_{\bu, \bv} \sum_{a \in \F_q}  \bra{\psi}(A^{\bv}_a \otimes I - I \otimes A^{\bv}_a)^2 \ket{\psi}}.
\end{align*}
The expression inside the first square root is at most~$1$ because~$T$ and~$H^{\bu}$ are sub-measurements,
and the expression inside the second square root is at most~$2\delta$ by \Cref{prop:simeq-to-approx} and the self-consistency of~$A$.
Now, \Cref{eq:move-over-v} is equal to
\begin{align*}
&\E_{\bu} \sum_{h, h'} \bra{\psi} H^{\bu}_{h'} \otimes (T_h\cdot \E_{\bv} A^{\bv}_{h(\bv)}) \ket{\psi}\\
=~& \E_{\bu} \sum_{h, h'} \bra{\psi} H^{\bu}_{h'} \otimes  (T_h\cdot Z)\ket{\psi} \tag{by \Cref{eq:swap-Z-for-A}}\\
=~& \E_{\bu} \sum_{ h'} \bra{\psi} H^{\bu}_{h'}  \otimes Z\ket{\psi} \tag{because~$T$ is a measurement}\\
\geq~& \E_{\bu} \sum_{ h'} \bra{\psi}  H^{\bu}_{h'} \otimes  \E_{\bv} A^{\bv}_{h'(\bv)} \ket{\psi} \tag{by \Cref{eq:Z-greater-than-A}}\\
=~& \E_{\bu, \bv} \sum_{ h'} \bra{\psi}  H^{\bu}_{h'} \otimes  A^{\bv}_{h'(\bv)} \ket{\psi}\\
=~& \E_{\bv} \sum_{a} \bra{\psi}  H_{[h(\bv)=a]} \otimes  A^{\bv}_{a} \ket{\psi}\\
=~& \E_{\bv} \sum_{a} \bra{\psi}  H_{[h(\bv)=a]} \otimes  I \ket{\psi} - \E_{\bv} \sum_{a\neq b} \bra{\psi}  H_{[h(\bv)=a]} \otimes  A^{\bv}_{b} \ket{\psi}\tag{because~$A$ is a measurement}\\
\geq~& \E_{\bv} \sum_{a} \bra{\psi}  H_{[h(\bv)=a]} \otimes  I \ket{\psi} - 4 \sqrt{\zeta_{\mathrm{variance}}} \tag{by \Cref{eq:explicit-bound-for-A-consistency}}\\
=~&  \sum_{h} \bra{\psi}  H_h \otimes  I \ket{\psi} - 4 \sqrt{\zeta_{\mathrm{variance}}}.
\end{align*}
In total, we have shown that
\begin{equation*}
\sum_{h\in\polyfunc{m}{q}{d}} \bra{\psi} H_h \otimes H_h \ket{\psi} \geq  \sum_{h} \bra{\psi}  H_h \otimes  I \ket{\psi} - 7 \sqrt{\zeta_{\mathrm{variance}}} - \sqrt{2\delta} - \frac{md}{q} - 4\sqrt{\zeta_{\mathrm{variance}}}.
\end{equation*}
Because
\begin{align*}
11 \sqrt{\zeta_{\mathrm{variance}}} + \sqrt{2\delta} + \frac{md}{q}
& = 11 \sqrt{24m\cdot\Big(\eps + \delta + \frac{md}{q}\Big)} + \sqrt{2\delta} + \frac{md}{q}\\
&\leq  55 m \cdot \Big(\eps^{1/2} + \delta^{1/2} + (d/q)^{1/2}\Big) + 2\sqrt{\delta} + m \cdot (d/q)^{1/2}\\
&\leq  57 m \cdot \Big(\eps^{1/2} + \delta^{1/2} + (d/q)^{1/2}\Big)\\
&\leq \zeta,
\end{align*}
this concludes the proof of \Cref{item:self-improvement-self}.

\vspace{\baselineskip}
\noindent
\emph{Proof of \Cref{item:self-improvement-boundedness} (Boundedness).}
The boundedness of~$H$ is

\begin{align*}
&\bra{\psi} Z \otimes I \ket{\psi} -\E_{\bu} \sum_a \bra{\psi}  A^{\bu}_{a} \otimes H_{[h(\bu)=a]} \ket{\psi} \\
=~&  \bra{\psi} Z \otimes I \ket{\psi} -\E_{\bu} \sum_h \bra{\psi}  A^{\bu}_{h(\bu)} \otimes H_h \ket{\psi}\\
=~&  \bra{\psi} Z \otimes I \ket{\psi} -\E_{\bu} \sum_h \bra{\psi}  I \otimes H_h \ket{\psi} +\E_{\bu} \sum_{a, h:h(\bu)\neq a} \bra{\psi}  A^{\bu}_{a} \otimes H_h \ket{\psi} \tag{because~$A$ is a measurement}\\
\leq~&  \bra{\psi} Z \otimes I \ket{\psi} -\E_{\bu} \sum_h \bra{\psi}  I \otimes H_h \ket{\psi} + 4 \sqrt{\zeta_{\mathrm{variance}}}\tag{by \Cref{eq:explicit-bound-for-A-consistency}}\\
\leq~& 3 \sqrt{\delta} + 4 \sqrt{\zeta_{\mathrm{variance}}}.\tag{by \Cref{eq:gonna-use-this-later-H-versus-Z}}
\end{align*}
We can bound the error by
\begin{align*}
 3 \sqrt{\delta} + 4 \sqrt{\zeta_{\mathrm{variance}}}
 & =  3 \sqrt{\delta} + 4 \sqrt{24m\cdot\Big(\eps + \delta + \frac{md}{q}\Big)}\\
 & \leq 3 \sqrt{\delta} + 20 m\cdot\Big(\eps^{1/2} + \delta^{1/2} + (d/q)^{1/2}\Big)\\
 &\leq 23 m\cdot\Big(\eps^{1/2} + \delta^{1/2} + (d/q)^{1/2}\Big)\\
 & \leq \zeta.
\end{align*}
This completes the proof.
\qed

\subsection{Self-improving to a projective measurement}
\label{sec:self-improvement-projective}

We now prove the full self-improvement theorem, i.e.\ \Cref{thm:self-improvement-in-induction-section}.
To do so, we will apply the orthonormalization lemma \Cref{thm:orthonormalization}
to the output of \Cref{lem:self-improvement-helper} and argue that it maintains the four properties of~$H$.

\begin{theorem}[Self-improvement; \Cref{thm:self-improvement-in-induction-section} restated]\label{thm:self-improvement}
Let $G \in \polymeas{m}{q}{d}$ be a measurement with the following properties:
\begin{enumerate}

  \item (Consistency with~$A$): \label{item:self-improvement-projective-G-consistency} On average over $\bu \sim \F_q^{m}$,
	  \begin{equation*}
	  A^{u}_a \otimes I \simeq_{\nu} I \otimes G_{[g(u)=a]}.
	  \end{equation*}
	  \end{enumerate}
Let
\begin{equation*}
\zeta = 3000m\cdot \Big(\eps^{1/32} + \delta^{1/32} + (d/q)^{1/32}\Big).
\end{equation*}
Then there exists a projective sub-measurement $H \in \polysub{m}{q}{d}$ with the following properties:
\begin{enumerate}
    \item(Completeness):  \label{item:self-improvement-projective-completeness} If $H =  \sum_h H_h$, then
    \begin{equation*}
    \bra{\psi} H \otimes I \ket{\psi} \geq (1-\nu)-\zeta.
    \end{equation*}
    \item(Consistency with~$A$):\label{item:self-improvement-projective-A-consistency} On average over $\bu \sim \F_q^m$,
        \begin{equation*}
            A^u_a \otimes I \simeq_{\zeta} I \otimes H_{[h(u) = a]}.
        \end{equation*}
    \item(Strong self-consistency): \label{item:self-improvement-projective-self}
    	\begin{equation*}
		H_h \otimes I \approx_{\zeta} I \otimes H_h.
	\end{equation*}
    \item(Boundedness): \label{item:self-improvement-projective-boundedness} There exists a positive-semidefinite matrix~$Z$ such that
    \begin{equation*}
        \bra{\psi} Z \otimes (I - H) \ket{\psi} \leq \zeta
    \end{equation*}
    and for each $h \in \polyfunc{m}{q}{d}$,
	\begin{equation*}
	Z \geq \left(\E_{\bu} A^{\bu}_{h(\bu)}\right).
	\end{equation*}
\end{enumerate}
\end{theorem}
\begin{proof}
We note that the bound we are proving is trivial when at least one of $\eps$, $\delta$, or $d/q$ is $\geq 1$.
In this case, $\zeta \geq 3000$.
Hence, we may assume that $\gamma, \zeta, d/q \leq 1$.
This will aid us when carrying out the error calculations near the end of the proof,
as it allows us to bound terms like $\eps^{1/2}$ by terms like $\eps^{1/4}$.

To begin,  apply \Cref{lem:self-improvement-helper} to~$G$.
Let $\widehat{H} \in \polysub{m}{q}{d}$ be the sub-measurement it outputs
and let $\widehat{\zeta}$ be the error
\begin{equation*}
\widehat{\zeta} = 100m \cdot \Big(\eps^{1/2} + \delta^{1/2} + (d/q)^{1/2}\Big).
\end{equation*}
By \Cref{item:self-improvement-self} of \Cref{lem:self-improvement-helper},
\begin{equation*}
\sum_{h} \bra{\psi} \widehat{H}_h \ot \widehat{H}_h \ket{\psi}
	\geq \bra{\psi} \widehat{H} \ot I \ket{\psi} - \widehat{\zeta}.
\end{equation*}
Thus, \Cref{thm:orthonormalization} implies the existence of a projective sub-measurement
$H \in \polysub{m}{q}{d}$ such that
\begin{equation}\label{eq:approx-between-H-with-and-without-hat}
\widehat{H}_h \otimes I \approx_{\widehat{\zeta}_{\mathrm{ortho}}} H_{h} \otimes I,
\end{equation}
where
\begin{equation*}
\widehat{\zeta}_{\mathrm{ortho}} = 100 \widehat{\zeta}^{1/4}.
\end{equation*}
In addition, \Cref{prop:self-consistency-implies-data-processing} implies that
\begin{equation}\label{eq:approx-data-processed}
\widehat{H}_{[h(u)=a]} \otimes I \approx_{\widehat{\zeta}_{\mathrm{dataprocess}}} H_{[h(u)=a]} \otimes I,
\end{equation}
where
\begin{equation*}
\widehat{\zeta}_{\mathrm{dataprocess}} = 8 \widehat{\zeta} + 8 \sqrt{\widehat{\zeta}_{\mathrm{ortho}}}.
\end{equation*}

Now we prove the four properties of this theorem.
We will show that each quantity is bounded by some error,
and then at the end of this proof we will show that all four errors are bounded by~$\zeta$.
\begin{enumerate}
    \item(Completeness):
\Cref{prop:completeness-transfer-self-consistent-A} implies that
\begin{align*}
 \bra{\psi} H \otimes I \ket{\psi}
 & \geq  \bra{\psi} \widehat{H} \otimes I \ket{\psi} - \widehat{\zeta} - 2\sqrt{\widehat{\zeta}_{\mathrm{ortho}}}\\
 & \geq  (1-\nu) -  2\widehat{\zeta} - 2\sqrt{\widehat{\zeta}_{\mathrm{ortho}}}. \tag{by \Cref{item:self-improvement-projective-completeness} of \Cref{lem:self-improvement-helper}}
\end{align*}
	\item (Consistency with~$A$):
\Cref{prop:triangle-sub} applied to
 \Cref{item:self-improvement-projective-A-consistency} of \Cref{lem:self-improvement-helper}
implies that
\begin{equation*}
A^u_a \otimes I \simeq_{\widehat{\zeta} + \sqrt{\widehat{\zeta}_{\mathrm{dataprocess}}}}
		I \otimes H_{[h(u) = a]}.
\end{equation*}
	\item (Strong self-consistency): \Cref{prop:two-notions-of-self-consistency} implies that
		\begin{equation*}
		\widehat{H}_h \otimes I \approx_{2\widehat{\zeta}} I \otimes \widehat{H}_h.
		\end{equation*}
		Thus,
		\begin{equation*}
			H_h \otimes I
			\approx_{\widehat{\zeta}_{\mathrm{ortho}}} \widehat{H}_{h} \otimes I
			\approx_{2\widehat{\zeta}} I \otimes \widehat{H}_{h}
			\approx_{\widehat{\zeta}_{\mathrm{ortho}}} I \otimes H_h.
		\end{equation*}
		Hence, by \Cref{prop:triangle-inequality-for-approx_delta},
		\begin{equation*}
		H_h \otimes I \approx_{6 \widehat{\zeta} + 6\widehat{\zeta}_{\mathrm{ortho}}} I \otimes H_h.
		\end{equation*}
	\item (Boundedness):
		The boundedness of~$H$ is
		\begin{align}
		&\bra{\psi} Z \otimes (I - H) \ket{\psi}\nonumber\\
		=~& \bra{\psi} Z \otimes I \ket{\psi} - \sum_h \bra{\psi} Z \otimes H_h \ket{\psi}\nonumber\\
		\leq~& \bra{\psi} Z \otimes I \ket{\psi} - \sum_h \bra{\psi} \E_{\bu} A^{\bu}_{h(\bu)} \otimes H_h \ket{\psi} \tag{by \Cref{item:self-improvement-projective-boundedness} of \Cref{lem:self-improvement-helper}}\\
		=~&  \bra{\psi} Z \otimes I \ket{\psi} -\E_{\bu} \sum_h \bra{\psi}  A^{\bu}_{h(\bu)} \otimes H_h \ket{\psi}\nonumber\\
		=~&  \bra{\psi} Z \otimes I \ket{\psi} -\E_{\bu} \sum_a \bra{\psi}  A^{\bu}_{a} \otimes H_{[h(\bu)=a]} \ket{\psi}. \label{eq:almost-there-self-improvement-edition}
		\end{align}
		By \Cref{prop:easy-approx-from-approx-delta},
		\begin{equation*}
		\E_{\bu} \sum_a \bra{\psi}  A^{\bu}_{a} \otimes H_{[h(\bu)=a]} \ket{\psi}
		\approx_{\sqrt{\widehat{\zeta}_{\mathrm{dataprocess}}}} \E_{\bu} \sum_a \bra{\psi}  A^{\bu}_{a} \otimes \widehat{H}_{[h(\bu)=a]} \ket{\psi}.
		\end{equation*}
		Hence,
		\begin{align*}
		\eqref{eq:almost-there-self-improvement-edition}
		& \leq \bra{\psi} Z \otimes I \ket{\psi} -\E_{\bu} \sum_a \bra{\psi}  A^{\bu}_{a} \otimes \widehat{H}_{[h(\bu)=a]} \ket{\psi} + \sqrt{\widehat{\zeta}_{\mathrm{dataprocess}}}\\
		& \leq \widehat{\zeta} + \sqrt{\widehat{\zeta}_{\mathrm{dataprocess}}}. \tag{by \Cref{item:self-improvement-boundedness} of \Cref{lem:self-improvement-helper}}
		\end{align*}
\end{enumerate}

This shows the four properties hold with errors
\begin{equation*}
2\widehat{\zeta} + 2\sqrt{\widehat{\zeta}_{\mathrm{ortho}}},\quad
\widehat{\zeta} + \sqrt{\widehat{\zeta}_{\mathrm{dataprocess}}},\quad
6 \widehat{\zeta} + 6\widehat{\zeta}_{\mathrm{ortho}},\quad
\text{and} \quad \widehat{\zeta} + \sqrt{\widehat{\zeta}_{\mathrm{dataprocess}}},
\end{equation*}
respectively. We now show that these four are bounded by~$\zeta$.
First, using the fact that $100^{1/4} \leq 4$, we note that
\begin{align*}
\widehat{\zeta}_{\mathrm{ortho}}
= 100 \widehat{\zeta}^{1/4}
&= 100 \Big(100m \cdot \Big(\eps^{1/2} + \delta^{1/2} + (d/q)^{1/2}\Big)\Big)^{1/4}\\
&\leq 400 m \cdot \Big(\eps^{1/8} + \delta^{1/8} + (d/q)^{1/8}\Big).
\end{align*}
Hence,
\begin{align*}
6 \widehat{\zeta} + 6\widehat{\zeta}_{\mathrm{ortho}}
& \leq 6 \Big(100m \cdot \Big(\eps^{1/2} + \delta^{1/2} + (d/q)^{1/2}\Big)\Big)
	+ 6 \Big(400 m \cdot \Big(\eps^{1/8} + \delta^{1/8} + (d/q)^{1/8}\Big)\Big)\\
&\leq 600m \cdot \Big(\eps^{1/8} + \delta^{1/8} + (d/q)^{1/8}\Big)
	+2400 m \cdot \Big(\eps^{1/8} + \delta^{1/8} + (d/q)^{1/8}\Big)\\
&= 3000m\cdot \Big(\eps^{1/8} + \delta^{1/8} + (d/q)^{1/8}\Big),
\end{align*}
which is less than~$\zeta$.
In addition, using the fact that~$\sqrt{400} = 20$, we also note that
\begin{align*}
\widehat{\zeta}_{\mathrm{dataprocess}}
&= 8 \widehat{\zeta} + 8 \sqrt{\widehat{\zeta}_{\mathrm{ortho}}}\\
&\leq 8 \Big(100m \cdot \Big(\eps^{1/2} + \delta^{1/2} + (d/q)^{1/2}\Big)\Big)
	+ 8 \sqrt{400 m \cdot \Big(\eps^{1/8} + \delta^{1/8} + (d/q)^{1/8}\Big)}\\
&\leq 800m \cdot \Big(\eps^{1/16} + \delta^{1/16} + (d/q)^{1/16}\Big)
	+ 160m \cdot \Big(\eps^{1/16} + \delta^{1/16} + (d/q)^{1/16}\Big)\\
&= 960m \cdot \Big(\eps^{1/16} + \delta^{1/16} + (d/q)^{1/16}\Big),
\end{align*}
which is clearly less than $\zeta$.
Thus, $2\widehat{\zeta} + 2\sqrt{\widehat{\zeta}_{\mathrm{ortho}}} \leq \widehat{\zeta}_{\mathrm{dataprocess}} \leq \zeta$.
Finally, using $\sqrt{960} \leq 31$,
\begin{align*}
\widehat{\zeta} + \sqrt{\widehat{\zeta}_{\mathrm{dataprocess}}}
&\leq  100m \cdot \Big(\eps^{1/2} + \delta^{1/2} + (d/q)^{1/2}\Big)
		+ \sqrt{960m \cdot \Big(\eps^{1/16} + \delta^{1/16} + (d/q)^{1/16}\Big)}\\
&\leq  100m \cdot \Big(\eps^{1/32} + \delta^{1/32} + (d/q)^{1/32}\Big)
		+ 31 m \cdot \Big(\eps^{1/32} + \delta^{1/32} + (d/q)^{1/32}\Big)\\
&=  131m \cdot \Big(\eps^{1/32} + \delta^{1/32} + (d/q)^{1/32}\Big),
\end{align*}
which is also less than~$\zeta$.
This completes the proof.
\end{proof}

